\documentclass[11pt]{article}
\usepackage[utf8]{inputenc}
\usepackage{amsmath,amssymb}
\usepackage[margin=1in]{geometry}
\usepackage{hyperref}
\emergencystretch=3em

\title{\texttt{LogBase} with a general tail rate:\\
$|G(L)-A|\le ML^{-\varepsilon}\ \Longrightarrow\ |H_m(L)-\textstyle\sum c_{m,i}\log^iL|\le (M/\varepsilon^m)L^{-\varepsilon}$}
\author{Claude, with Daniel Murfet}
\date{2026-09-06}

\begin{document}
\maketitle
\sloppy

\begin{abstract}
Unit~49's base predicate required the tail $|G(L)-A|\le M/L$. The mixed-multiplicity case needs bases
whose tail decays only like $L^{-(q-p)}$, where $q>p$ are the non-minimal and minimal candidate
exponents. This unit generalises \texttt{LogBase} by a rate $\varepsilon>0$ and shows that the
transfer step preserves the rate at the cost of a factor $1/\varepsilon$ per level:
$|\theta_m(L)|\le(M/\varepsilon^m)L^{-\varepsilon}$. The weighted primitives are the case
$\varepsilon=1$, and every downstream statement is unchanged. Zero \texttt{sorry}/\texttt{axiom}.
\end{abstract}

\section{Setting}
With $|\theta_m(x)|\le M'x^{-\varepsilon}$ the tail integrand satisfies $|\theta_m(x)/x|\le
M'x^{-\varepsilon-1}$, integrable on $[1,\infty)$ for any $\varepsilon>0$, and
$\int_L^\infty M'x^{-\varepsilon-1}=M'L^{-\varepsilon}/\varepsilon$. So the recursion
$\theta_{m+1}(L)=-\int_L^\infty\theta_m/x$ yields $|\theta_{m+1}(L)|\le(M'/\varepsilon)L^{-\varepsilon}$.

\section{The things formalised}
{\small
\begin{verbatim}
structure LogBase (G : ℝ → ℝ) (A M C δ ε : ℝ) : Prop where
  ...
  ε_pos : 0 < ε
  tail : ∀ L, 1 ≤ L → |G L - A| ≤ M * L ^ (-ε)
theorem divRem_tail_abs_le (G : ℝ → ℝ) (A M ε : ℝ) (hε : 0 < ε) (m : ℕ) (L : ℝ) (hL : 1 ≤ L)
    (hb : ∀ L, 1 ≤ L → |divRem G A m L| ≤ M * L ^ (-ε)) :
    |∫ x in Ioi L, divRem G A m x / x| ≤ M / ε * L ^ (-ε)
theorem divRem_abs_le (hG : LogBase G A M C δ ε) (m : ℕ) :
    ∀ L, 1 ≤ L → |divRem G A m L| ≤ M / ε ^ m * L ^ (-ε)
theorem iterDivPrim_expansion (hG : LogBase G A M C δ ε) (m : ℕ) (L : ℝ) (hL : 1 ≤ L) :
    |iterDivPrim G m L - ∑ i ∈ Finset.range (m + 1), divCoeff G A m i * Real.log L ^ i|
      ≤ M / ε ^ m * L ^ (-ε)
\end{verbatim}
}

\section{Proof strategy}
A mechanical generalisation: the exponent $-2$ becomes $-\varepsilon-1$ in the domination and in
\texttt{integral\_Ioi\_rpow\_of\_lt}, the pointwise bound uses \texttt{Real.rpow\_sub\_one}, and the
induction carries $M/\varepsilon^m$ via \texttt{pow\_succ} and \texttt{div\_div}. In the asymptotic of
unit~50 the remainder is bounded by $M/\varepsilon^m$ on $L\ge1$ using
\texttt{Real.rpow\_le\_one\_of\_one\_le\_of\_nonpos}. The $F_\gamma$ instance sets $\varepsilon=1$ and
rewrites $L^{-1}$ as $1/L$ so the previously stated corollary keeps its form.

\section{Roadblocks and resolutions}
\begin{itemize}
\item None: both files compiled on the first pass after the edit. The only care was to thread the
level-dependent constant $M'$ through the lemmas that take the remainder bound as a hypothesis.
\end{itemize}

\section{What was Mathlib, what was new}
Mathlib: \texttt{Real.rpow\_sub\_one}, \texttt{integral\_Ioi\_rpow\_of\_lt}, and the bound
\texttt{Real.rpow\_le\_one\_of\_one\_le\_of\_nonpos} for $L^{-\varepsilon}\le1$. New: the rate
parameter and its propagation.

\section{Lessons}
When an abstraction is introduced (unit~49), the first consumer beyond the motivating instance usually
exposes a hypothesis that was accidentally too strong; here the fix was a five-minute generalisation
because the abstraction had isolated it.

\section{Outcome}
The iterated expansion is available for bases with any polynomial tail rate. Next: the mixed case
with one noncritical coordinate innermost, whose base $G_1(y)=\int_0^yx^{p-q-1}F_{q-1}$ is a
\texttt{LogBase} with $A=C(p,q)$, $\varepsilon=q-p$, giving $Z\sim c\,n^{-p/2}(\log n)^{m-1}$ with
$m$ the number of minimal coordinates.
\end{document}
